% Native reconstruction of Springer SGA4 tome II backmatter.
% Source SHA-256 7c5b8cef5a99d22785501eead8353641064f858c51ce361b00b310b2bcc46585.
% Physical pages415--420 (zero-based414--419), printed pages413--418.
% All entries retain original order and every printed numerical reference.
\chapter*{术语索引}
\begingroup
\small
\setlength{\LTleft}{0pt}
\setlength{\LTright}{0pt}
\begin{longtable}{@{}p{.68\linewidth}p{.27\linewidth}@{}}
% SOURCE physical415, printed413: all23 entries directly checked against scan.
% Algébrique (topos, objet)
代数的 (拓扑斯、对象) & VI 2.3\\
% Augmentation de descente effective
有效下降增广 & $\mathrm{V}^{\mathrm{bis}}$ 2.2.2\\
% Biscindage
双分裂 & VI 7.1.4\\
% Cartésien (morphisme)
笛卡儿 (态射) & VI 6.1\\
% Cartésiens (foncteurs)
笛卡儿 (函子) & VI 6.1.4\\
% Catégorie localement filtrante
局部滤过范畴 & V 8.1.0\\
% Centre (d'un topos local)
(局部拓扑斯的) 中心 & VI 8.4.6\\
% Cohérent (morphisme)
凝聚 (态射) & VI 1.7\\
& VI 1.21\\
% Cohérent (morphisme de topos)
凝聚 (拓扑斯态射) & VI 3.1\\
& VI 8.3.14\\
& VI 3.3\\
% Cohérent (topos)
凝聚 (拓扑斯) & VI 8.3.13\\
& VI 2.3\\
& VI 5.2\\
% Cohérent (topos localement -)
局部凝聚 (拓扑斯) & VI 2.3\\
% Cohomologie à support dans Z (groupe de, faisceau de)
支集包含于 $Z$ 的上同调 (群、层) & V 6.3\\
% Cohomologie de Čech (groupe)
Čech 上同调 (群) & V 2.4.3\\
% Cohomologie de Čech (préfaisceau)
Čech 上同调 (预层) & V 2.4.5\\
% Complexe de Čech
Čech 复形 & V 2.4.3\\
% Constructible
可构造 & VI 1.9.3\\
& VI 1.21\\
% Descente cohomologique effective
有效上同调下降 & $\mathrm{V}^{\mathrm{bis}}$ 2.2.6\\
% Descente cohomologique relative
相对上同调下降 & $\mathrm{V}^{\mathrm{bis}}$ 2.4\\
% G-1-descente cohomologique
$G$-$1$-上同调下降 & $\mathrm{V}^{\mathrm{bis}}$ 3.1.2\\
% Faisceaux acycliques. No reference is printed in the source.
无环层 & \\
% Faisceaux C-acycliques
$C$-无环层 & V 4.1\\
% Faisceaux Ext à supports dans Z
支集包含于 $Z$ 的 $\operatorname{Ext}$ 层 & V 6.8.2\\
% Faisceaux flasques
松弛层 & V 4.1\\[.7em]
% SOURCE physical416, printed414: all28 entries directly checked against scan.
% Famille de supports
支集族 & V 6.12\\
% Famille de supports de caractère local
具有局部特征的支集族 & V 6.12.2\\
% Fibrée (catégorie)
纤维化 (范畴) & VI 6.1\\
% Fibre (catégorie)
纤维 (范畴) & VI 6.1.1\\
% Foncteur cofinal (entre catégories localement filtrantes)
(局部滤过范畴之间的) 共尾函子 & V 8.1.4\\
% Foncteur fibre
纤维函子 & VIII 3.3\\
% Groupe Ext à support dans Z
支集包含于 $Z$ 的 $\operatorname{Ext}$ 群 & V 6.8.5\\
% Homologie
同调 & V 7.2.0\\
% Homotopie
同伦 & V 7.3.1\\
% Hyper-recouvrement
超覆盖 & V 7.3.1\\
% Limite inductive (d'une catégorie fibrée)
(纤维范畴的) 归纳极限 & VI 6.3\\
% Limite inductive locale
局部归纳极限 & V 8.2.1\\
% Limite projective (d'une catégorie fibrée)
(纤维范畴的) 投射极限 & VI 6.10\\
% Limite projective de topos fibrés
纤维化拓扑斯的投射极限 & VI 8.1.1\\
% Localisé strict
严格局部化 & VIII 4.3\\
% Morphisme de U-sites fibrés
纤维化 $\underline U$-景的态射 & VI 7.2.2\\
% Morphismes cartésiens de U-sites fibrés
纤维化 $\underline U$-景的笛卡儿态射 & VI 7.2.2\\
% Morphisme spécial (de préfaisceaux simpliciaux)
(单纯预层的) 特殊态射 & V 7.3.3\\
% Noethérien
Noether 的 & VI 2.11\\
% Objet semi-simplicial tronqué
截断半单纯对象 & V 7.1.0\\
% Parfait (topos)
完美 (拓扑斯) & VI 1.27\\
% Plats (Modules)
平坦 (模) & V 1.1\\
% Plat (Morphisme de topos)
平坦 (拓扑斯态射) & V 1.8\\
% Point géométrique
几何点 & VIII 3.1\\
% Prénoethérien
预 Noether 的 & VI 1.30\\
% Quasi-compact (objet)
拟紧 (对象) & VI 1.1\\
& VI 1.21\\
% Quasi-compact (morphisme)
拟紧 (态射) & VI 1.7\\
& VI 1.21\\
% Quasi-compact (morphisme de topos)
拟紧 (拓扑斯态射) & VI 3.1\\[.7em]
% SOURCE physical417, printed415: all28 entries directly checked against scan.
% Quasi-séparé (morphisme)
拟分离 (态射) & VI 1.7\\
& VI 1.21\\
% Quasi-séparé (morphisme de topos)
拟分离 (拓扑斯态射) & VI 3.1\\
% Quasi-séparé (topos)
拟分离 (拓扑斯) & VI 2.3\\
% Résolution Standard
标准消解 & V 2.3.6\\
% Sections cartésiennes
笛卡儿截面 & VI 6.10\\
% Semi-représentable
半可表 & V 7.3\\
% U-site fibré
纤维化 $\underline U$-景 & VI 7.2.1\\
% Site total
全景 & VI 7.4.1\\
% Squelette, Cosquelette
骨架、余骨架 & V 7.1.1\\
% Strictement local (anneau, schéma)
严格局部 (环、概形) & VIII 4.2\\
% Suite spectrale de Cartan-Leray (relative à un morphisme de topos)
(关于拓扑斯态射的) Cartan-Leray 谱序列 & V 5.3\\
% Suite spectrale de Cartan-Leray (relative à un recouvrement)
(关于覆盖的) Cartan-Leray 谱序列 & V 3.2\\
% Suite spectrale de descente
下降谱序列 & VIII 8.1\\
% Topologie de Zariski
Zariski 拓扑 & VII 4.2\\
% Topologie étale
平展拓扑 & VII 1.2\\
% Topologie fidèlement plate quasi-compacte
忠实平坦拟紧拓扑 & VII 4.2\\
% Topologie totale (d'un U-site fibré)
(纤维化 $\underline U$-景的) 全拓扑 & VI 7.4.1\\
% D-topos
$D$-拓扑斯 & $\mathrm{V}^{\mathrm{bis}}$ 1.2.1\\
% D-topos annelé
赋环 $D$-拓扑斯 & $\mathrm{V}^{\mathrm{bis}}$ 1.3.1\\
% Topos classifiant
分类拓扑斯 & VIII 2.1\\
% Topos étale
平展拓扑斯 & VIII 1.2\\
% Topos fibré
纤维化拓扑斯 & VI 7.1\\
% Topos fibré annelé
赋环纤维化拓扑斯 & VI 8.6.1\\
% Topos fibré annelé plat
平坦赋环纤维化拓扑斯 & VI 8.6.4\\
% U-topos fibré associé
相伴纤维化 $\underline U$-拓扑斯 & VI 7.2.6\\
% Topos fini, profini
有限、pro-有限拓扑斯 & VI 2.9.4\\
% Topos local
局部拓扑斯 & VI 8.4.6\\
% Topos localisé
局部化拓扑斯 & VI 8.4.2\\[.7em]
% SOURCE physical418, printed416: both entries directly checked against scan.
% Topos simplicial
单纯拓扑斯 & $\mathrm{V}^{\mathrm{bis}}$ 1.2.1\\
% Topos total
全拓扑斯 & VI 7.4.3.3\\
\end{longtable}
\endgroup

\chapter*{记号索引}
\begingroup
\small
\setlength{\LTleft}{0pt}
\setlength{\LTright}{0pt}
\begin{longtable}{@{}p{.68\linewidth}p{.27\linewidth}@{}}
% SOURCE physical419, printed417, left column: all27 entries in column order.
$A^{\mathrm{hs}}$ & VIII 4.4\\
$A_{\bullet}(\mathcal U)$ & V 1.10.3\\
$C^{\sim/I}$ & VI 7.2.6.6\\
$C^{\bullet}(\mathcal U,M)$ & V 2.3.3.1\\
$E_{\mathrm{coh}},E_{\mathrm{PF}}$ & VI 1.2.4\\
& VI 8.3.13\\
$E_{\mathrm{cons}}$ & VI 2.9.3\\
$\operatorname{Et}_{/X}$ & VII 1.3\\
$\operatorname{Ext}^{q}_{A}(E,M,N)$ & V 2.1.1\\
$\operatorname{Ext}^{q}_{A,Z}(X;M,N)$ & V 6.8.5\\
$\operatorname{Ext}^{q}_{A,\Phi}(X;\mathord\cdot,\mathord\cdot)$ & V 6.13.8\\
$\mathcal{E}\!xt^{q}_{A}(M,N)$ & V 6.0.1\\
$\mathcal{E}\!xt^{q}_{A,Z}(M,N)$ & V 6.8.2\\
$\mathcal{E}\!xt^{q}_{A,\Phi}(\mathord\cdot,\mathord\cdot)$ & V 6.13.8\\
$F\mapsto F_{\xi}$ & VIII 3.3\\
$\underrightarrow F,(\underrightarrow F)^N$ & VI 8.2.3\\
% The source visibly omits both morphism names before these colons.
${}:F\to F_i,\quad{}:F\times I\to F$ & VI 8.1.3\\
$f^*,f_*$ & VII 1.4\\
$(\mathrm{fpqc})$ & VII 4.2\\
$\mathcal F(S^{-1})$ & VI 6.2\\
& VI 6.5\\
$H^*(\pi,M)$ & VIII 2.3\\
$H^i(X,F)$ & VII 1.3\\
$\operatorname{HR}_p$ & V 7.3\\
$\underline{\operatorname{HR}}_p$ & V 7.3.1.6\\
$\Hvr{q},\Hvv{q}$ & V 7.4.0\\
$H^q(E,N)$ & V 2.1.1\\
$H^q(\mathcal U,M)$ & V 2.3.3\\
$H^q(X,M)$ & V 2.4.1\\[.7em]
% SOURCE physical419, printed417, right column: all27 entries in column order.
$H^q(X,N)$ & V 2.1.1\\
$H^q_{\Phi}$ & V 6.13.8\\
$\underline H^{q}_{\Phi}$ & V 6.13.8\\
$H^q_Z(E,\mathord\cdot)$ & V 6.3\\
$\underline H^{q}_Z(\mathord\cdot)$ & V 6.3\\
$\mathcal H^{q}(M)$ & V 2.4.2.2\\
$\check{\mathcal H}^{q}(G)$ & V 2.4.5.2\\
$\underline{\operatorname{Hom}}_{\mathrm{Cart}/I}(F,G)$ & VI 6.1.4\\
& VI 6.9\\
$\operatorname{Hom}_f(X,Y)$ & VI 6.0\\
$\underline{\operatorname{Hom}}_I(F,G)$ & VI 6.0\\
$\operatorname{Hom}_{\xi}(X,Y)$ & VI 6.1.1\\
$\underline{\operatorname{Homtop}}_{\mathrm{Cart}/I}(F,G)$ & VI 7.1.7\\
$\underline{\operatorname{Homtop}}_I(F,G)$ & VI 7.1.7\\
${\varprojlim}^{q}$ & V 2.3.1\\
$\underrightarrow{\operatorname{Lim}}_{\mathcal E^{\circ}}$ & VI 6.3\\
$\underleftarrow{\operatorname{Lim}}_{\mathcal E^{\circ}}\mathcal F$ & VI 6.9\\
$\underleftarrow{\operatorname{Limtop}}_I F,\quad\underleftarrow{\operatorname{Limtop}}_I F_i,\quad\underleftarrow F$ & VI 8.1.3\\
$\operatorname{Loc}_p(E)$ & VI 8.4.2.1\\
$\underline{\operatorname{Morsite}}_{\mathrm{Cart}/I}(C,D)$ & VI 7.2.2.2\\
$\underline{\operatorname{Morsite}}_I(C,D)$ & VI 7.2.2\\
$Q:F\to\operatorname{Top}(F)$ & VI 8.5.1.2\\
$R^q u_*$ & V 5.0\\
$[R|_fS],[R|S]$ & $\mathrm{V}^{\mathrm{bis}}$ 1.2.7\\
$(\operatorname{Sch})$ & VII 1.1\\
$\operatorname{sk}_n,\operatorname{cosk}_n$ & V 7.1.1\\
$S_{\bullet}(\mathcal U)$ & V 1.10\\
$\operatorname{Top}(C)$ & VI 7.4.3.3\\[.7em]
% SOURCE physical420, printed418: all7 entries directly checked against scan.
$X_{\text{ét}}$ & VII 1.2\\
$\widetilde X_{\text{ét}}$ & VII 1.2\\
$[[X|_uX']]$ & $\mathrm{V}^{\mathrm{bis}}$ 1.2.7\\
$X_{\mathrm{Zar}}$ & VII 4.2.2\\
$\alpha_i:F_i\to F$ & VI 7.4.3.4\\
& VI 7.4.12\\
$\Delta,\Delta(E),\Delta[n],\Delta E[n],\Delta_n$ & V 7.1.0\\
$\varepsilon_{C/I}:C\to C^{\sim/I}$ & VI 7.2.6.7\\
& VI 7.4.11\\
\end{longtable}
\endgroup
